% ============================================================================
% Section VIII: Conclusions and Future Work
% ============================================================================

\section{Conclusions}
\label{sec:conclusions}

A novel, high-performance tool for Generation and Transmission Expansion
Planning, designated \gtopt{}, has been presented.  The main conclusions
are as follows:

\begin{enumerate}
  \item \textbf{Complete formulation}: The GTEP problem is formulated
    as a sparse LP/MIP covering DC optimal power flow with Kirchhoff
    constraints, battery energy storage with state-of-charge tracking
    and self-discharge losses, hydro cascade systems with reservoirs,
    turbines, filtration, and volume-dependent efficiency, a generic
    water rights layer for irrigation agreement modeling with
    volume-dependent dynamic bounds, spinning reserve requirements,
    and multi-stage capacity expansion with annualized investment
    costs.

  \item \textbf{Three solution methods}: The monolithic LP provides
    exact solutions for moderate-size problems.  SDDP with
    multi-scenario apertures handles large stochastic problems
    efficiently, converging to within 1\% of the monolithic optimum
    in 10--15 iterations.  The novel cascade decomposition combines
    both approaches for long-horizon planning.

  \item \textbf{High performance}: The C++26 implementation with
    flat-map-based sparse-matrix assembly achieves a 10--27$\times$
    speedup over Python-based alternatives in LP assembly.  For the
    largest configuration tested (57-bus, 8760 blocks, ${\sim}18$M
    variables), assembly completes in approximately 8 seconds.
    Pluggable solver backends (CLP, CBC, CPLEX, HiGHS) are supported
    through a dynamic plugin architecture.

  \item \textbf{Validated correctness}: Cross-validation against
    pandapower DC OPF on IEEE benchmark networks (4--57 buses)
    demonstrates agreement to solver tolerance ($<10^{-6}$ relative
    error).  All 13 benchmark cases converge to optimal solutions.

  \item \textbf{Comprehensive toolchain}: A Python companion toolchain
    (igtopt, plp2gtopt, pp2gtopt, gtopt\_compare) and web interfaces
    provide a complete workflow from data preparation through result
    analysis, supporting interoperability with existing planning
    ecosystems.
\end{enumerate}

\subsection{Future Work}

Planned enhancements include:

\begin{itemize}
  \item \textbf{AC OPF support}: Extension of the formulation to include
    reactive power and voltage magnitude constraints via second-order
    cone programming (SOCP) relaxations.
  \item \textbf{Unit commitment}: Addition of binary start-up/shut-down
    decisions and minimum up/down time constraints for thermal generators.
  \item \textbf{Multi-sector coupling}: Extension of the model to include
    gas networks, district heating, and hydrogen production/storage for
    integrated energy system planning.
  \item \textbf{GPU-accelerated LP assembly}: Exploitation of GPU
    parallelism for sparse-matrix construction in very large-scale
    problems.
  \item \textbf{Machine learning integration}: Application of trained
    models to warm-start the SDDP algorithm or provide scenario
    reduction for stochastic programming.
  \item \textbf{Real-world case studies}: Application of \gtopt{} to
    national-scale expansion planning problems with detailed hydro
    cascade modeling.
\end{itemize}

The \gtopt{} tool is designed to support reproducible research in power
system optimization and to facilitate community collaboration in future
releases.
